\section{Interpretation and further analytical questions}
\label{sec:discussion}

The counted quadratic factor is the common starting point of the paper.
Its symmetric coordinates remain regular at a merger, and its rational
traces remove the need to follow individual exterior roots. Divided
differences then give the complete integer hierarchy. Split-pole moments
are algebraically equivalent coordinates, while the complex exponent
organizes a holomorphic family and complete damped contour formulas.
Neither change of representation alone supplies the signed information
needed by gap transport.

The recent consequences sharpen that distinction in two directions.
The finite-head residual is explicitly small in moment order, yet its
independent positive comparison can grow factorially. The appropriate
question is therefore whether compatibility or a finite-divisor
reconstruction yields a useful error theorem. Separately, the joint
forcing, weighted gap energy and deflated derivative expansion retain
more coupling than separate absolute bounds. Their exactness identifies
which additional estimate is needed, but does not supply it. The scalar
energy example concerns the two-equation consumer only, and the static
divisor example concerns the stated enlarged-constant slice class.
Neither excludes a better estimate from the full canonical dynamics.

In the all-real regime there is an additional elementary sign structure:
the complete exterior denominators of an adjacent pair are positive.
The backward comparison uses it to improve the local safety allowance
from initial neighboring gaps. Regional minima make extremal selection
well-defined without a global minimizing pair, and their weak evolution
handles switches of the minimizing label. The distribution nevertheless
couples to the exterior interaction. A joint law of consecutive gaps is
therefore a concrete statistical input, with a computable consequence,
rather than a substitute name for the desired collision exclusion.

Two analytical tasks remain distinct. First, a finite canonical product
could provide a compatible approximation on a common interval if its
complete error, local contour clearance and moment transfer are proved.
The polynomial identity in \eqref{fin:polynomial-heat} and the existing
slice-jet passage are useful ingredients, not the missing time-uniform
theorem. Second, a statistical descent requires an actual estimate on the
lower tail of the three-gap safety margin, together with an argument that
eliminates all exceptions over the interval and at unbounded heights.
A density-one conclusion or fixed positive-blur modulus alone does not
meet that requirement. For an isolated very narrow gap, the neighboring
comparison gives little relative improvement, making those configurations
a substantive part of the remaining problem.

The statistical use of gap evolution has an important precedent in
Rodgers and Tao \cite{RodgersTao2020Newman}. Their contradiction under a
negative Newman constant gives the lower bound $\Lambda\ge0$. The
proposed application here has the other direction: continue an established
all-real time downward to improve an upper bound. The local comparison
does not provide the complete initial statistics or uniform height
coverage needed for that conclusion. No improved value of $\Lambda$
follows from the present results.

The collision controls in the project background are consistent with
these scopes. Quadratic and quartic heat polynomials obey the local
algebra while lacking the canonical growing population. Source
perturbations with Euler translates, finite square weights or finite
prime constraints change the full reflected unit-weight theta source;
the documented endpoint-defective subclasses also fail the fixed strip.
Reflected super-Gaussian controls fail the order-one product premise.
Smooth positive and positive-primitive controls alter theta arithmetic
and do not establish the complete strip/count/safe-line inputs used
here. Atomic or finite-profile controls lack the full canonical
population or profile; changed observations and fractional profiles
change the physical readout or integer multiplicities. RH-conditional
endpoint controls retain RH as an input and change ordinary Dirichlet
arithmetic. The static surgery of Proposition~\ref{prop:static}
has no common canonical heat tube, and the scalar family of
Proposition~\ref{joint:scalar-screen} has no asserted divisor realization.
The detailed source and premise screen is retained in
\cite{BakerControls,BakerCMReview,BakerCM13,BakerCM17}.

All examples used for the conclusions of this draft are analytical.
The quartic controls are exact heat polynomials, the static construction
is an asymptotic entire-function family, and the truncation limitation
is an explicit path in a positive matrix exponential. Numerical work may
support a future specified estimate, but the proofs here do not depend
on samples or floating-point evidence. The unresolved questions concern
complete-divisor, common-time and uniform-height control.
